\documentclass{beamer}

\usepackage{amsmath,amsfonts,amssymb,bm}
\usepackage{quantikz}
\usepackage{braket}
\usepackage{graphicx}

\usetheme{Madrid}

\title{Quantum Support Vector Machines}
\subtitle{Kernel Methods, Quantum Feature Maps, and Circuits}
\author{}
\date{}

\begin{document}

\frame{\titlepage}

%================================================
\section{Motivation}
%================================================

\begin{frame}{Why Quantum Machine Learning?}
\begin{itemize}
\item Classical ML struggles with high-dimensional feature spaces
\item Kernel methods implicitly lift data into large Hilbert spaces
\item Quantum systems naturally live in exponentially large Hilbert spaces
\end{itemize}

\medskip

\begin{block}{Key Idea}
Use a quantum computer to efficiently compute inner products in large Hilbert spaces.
\end{block}
\end{frame}

%------------------------------------------------

\begin{frame}{Physics Viewpoint}
\begin{itemize}
\item Data $\rightarrow$ quantum state $\ket{\phi(x)}$
\item Kernel:
\[
K(x,x') = |\braket{\phi(x)}{\phi(x')}|^2
\]
\item This is a transition probability between quantum states
\end{itemize}

\medskip

Thus QSVMs are naturally formulated in quantum mechanics.
\end{frame}

%================================================
\section{Classical Support Vector Machines}
%================================================

\begin{frame}{Binary Classification}
Given labeled data:
\[
(x_i, y_i), \quad y_i \in \{-1,1\}
\]

Find hyperplane:
\[
w \cdot x + b = 0
\]
\end{frame}

%------------------------------------------------

\begin{frame}{Maximum Margin Principle}
Solve:
\[
\min_{w,b} \frac{1}{2}\|w\|^2
\]
subject to:
\[
y_i(w\cdot x_i + b) \ge 1
\]
\end{frame}

%------------------------------------------------

\begin{frame}{Dual Formulation}
\[
\max_{\alpha}
\sum_i \alpha_i
-
\frac{1}{2}
\sum_{i,j} \alpha_i \alpha_j y_i y_j K(x_i,x_j)
\]

with:
\[
K(x_i,x_j) = x_i \cdot x_j
\]
\end{frame}

%------------------------------------------------

\begin{frame}{Kernel Trick}
Replace inner product:
\[
x_i \cdot x_j \rightarrow K(x_i,x_j)
\]

\medskip

Implicit mapping:
\[
x \rightarrow \phi(x)
\]

\[
K(x,x') = \phi(x) \cdot \phi(x')
\]
\end{frame}

%================================================
\section{Quantum Feature Maps}
%================================================

\begin{frame}{Quantum Embedding}
Encode classical data into quantum states:
\[
\ket{\phi(x)} = U(x)\ket{0}
\]
\end{frame}

%------------------------------------------------

\begin{frame}{Quantum Kernel}
\[
K(x,x') = |\braket{\phi(x)}{\phi(x')}|^2
\]

\begin{itemize}
\item computed via quantum circuits
\item corresponds to state overlap
\end{itemize}
\end{frame}

%------------------------------------------------

\begin{frame}{Example Feature Map}
Single qubit:
\[
U(x) = e^{i x Z}
\]

\[
\ket{\phi(x)} = \cos x \ket{0} + i \sin x \ket{1}
\]
\end{frame}

%------------------------------------------------

\begin{frame}{Multi-Qubit Feature Map}
\[
U(x) = \prod_i e^{i x_i Z_i} \prod_{i<j} e^{i x_i x_j Z_i Z_j}
\]

\begin{itemize}
\item introduces entanglement
\item encodes nonlinear correlations
\end{itemize}
\end{frame}

%================================================
\section{Quantum Kernel Estimation}
%================================================

\begin{frame}{Overlap Computation}
We need:
\[
|\braket{\phi(x)}{\phi(x')}|^2
\]

Two main approaches:
\begin{itemize}
\item swap test
\item direct fidelity estimation
\end{itemize}
\end{frame}

%------------------------------------------------

\begin{frame}{Swap Test Circuit}
\[
\begin{quantikz}
\lstick{\ket{0}} & \gate{H} & \ctrl{1} & \gate{H} & \meter{} \\
\lstick{\ket{\phi(x)}} & \qw & \swap{1} & \qw & \qw \\
\lstick{\ket{\phi(x')}} & \qw & \targX{} & \qw & \qw
\end{quantikz}
\]

Probability of measuring $0$:
\[
P(0) = \frac{1 + |\braket{\phi(x)}{\phi(x')}|^2}{2}
\]
\end{frame}

%------------------------------------------------

\begin{frame}{Kernel Matrix}
Construct:
\[
K_{ij} = |\braket{\phi(x_i)}{\phi(x_j)}|^2
\]

\begin{itemize}
\item evaluated on quantum hardware
\item used in classical SVM solver
\end{itemize}
\end{frame}

%================================================
\section{QSVM Algorithm}
%================================================

\begin{frame}{Hybrid QSVM}
\begin{enumerate}
\item Encode data into quantum states
\item Compute kernel matrix using quantum circuits
\item Solve SVM optimization classically
\end{enumerate}
\end{frame}

%------------------------------------------------

\begin{frame}{Decision Function}
\[
f(x) = \sum_i \alpha_i y_i K(x_i,x) + b
\]
\end{frame}

%================================================
\section{Example Circuits}
%================================================

\begin{frame}{Feature Map Circuit}
\[
\begin{quantikz}
\lstick{\ket{0}} & \gate{R_z(x_1)} & \ctrl{1} & \qw \\
\lstick{\ket{0}} & \gate{R_z(x_2)} & \targ{} & \qw
\end{quantikz}
\]

\begin{itemize}
\item encodes classical data
\item introduces entanglement
\end{itemize}
\end{frame}

%------------------------------------------------

\begin{frame}{Kernel Evaluation Circuit}
\[
\begin{quantikz}
\lstick{\ket{0}} & \gate{U(x)} & \gate{U^\dagger(x')} & \meter{}
\end{quantikz}
\]

\[
P(0) = |\braket{\phi(x)}{\phi(x')}|^2
\]
\end{frame}

%================================================
\section{Physics Interpretation}
%================================================

\begin{frame}{Hilbert Space Perspective}
\begin{itemize}
\item classical SVM: feature space $\mathbb{R}^d$
\item QSVM: quantum Hilbert space $\mathcal{H}$
\end{itemize}

\medskip

\[
\dim(\mathcal{H}) = 2^n
\]
\end{frame}

%------------------------------------------------

\begin{frame}{Connection to Many-Body Physics}
\begin{itemize}
\item feature maps $\leftrightarrow$ quantum states
\item kernels $\leftrightarrow$ overlaps
\item classification $\leftrightarrow$ phase discrimination
\end{itemize}
\end{frame}

%------------------------------------------------

\begin{frame}{Expressivity}
\begin{itemize}
\item entangled feature maps generate complex kernels
\item difficult to simulate classically
\end{itemize}
\end{frame}

%================================================
\section{Advantages and Limitations}
%================================================

\begin{frame}{Advantages}
\begin{itemize}
\item access to exponential feature space
\item potentially classically hard kernels
\end{itemize}
\end{frame}

%------------------------------------------------

\begin{frame}{Challenges}
\begin{itemize}
\item noise and decoherence
\item kernel estimation cost
\item unclear quantum advantage in practice
\end{itemize}
\end{frame}

%================================================
\section{Extensions}
%================================================

\begin{frame}{Variational QSVM}
\begin{itemize}
\item trainable feature maps
\item hybrid optimization
\end{itemize}
\end{frame}

%------------------------------------------------

\begin{frame}{Connection to Other Methods}
\begin{itemize}
\item quantum neural networks
\item quantum kernel methods
\item variational classifiers
\end{itemize}
\end{frame}

%================================================
\section{Summary}
%================================================

\begin{frame}{Summary}
\begin{itemize}
\item QSVM = kernel SVM with quantum feature map
\item kernel:
\[
K(x,x') = |\braket{\phi(x)}{\phi(x')}|^2
\]
\item computed via quantum circuits
\item connects ML with quantum mechanics
\end{itemize}
\end{frame}


%================================================
\section{Pauli Expansion of Feature Maps}
%================================================

\begin{frame}{General Quantum Feature Map}
A quantum feature map is a unitary embedding:
\[
\ket{\phi(x)} = U(x)\ket{0}
\]

We can express the unitary as
\[
U(x) = e^{i H(x)},
\]
where \(H(x)\) is a Hermitian operator depending on the data.
\end{frame}

%------------------------------------------------

\begin{frame}{Pauli Decomposition}
Any Hermitian operator can be expanded as
\[
H(x) = \sum_{\alpha} f_\alpha(x) P_\alpha,
\]
where:
\begin{itemize}
\item \(P_\alpha \in \{I,X,Y,Z\}^{\otimes n}\),
\item \(f_\alpha(x)\) are real-valued functions.
\end{itemize}

Thus:
\[
U(x) = e^{i \sum_\alpha f_\alpha(x) P_\alpha}.
\]
\end{frame}

%------------------------------------------------

\begin{frame}{Example: ZZ Feature Map}
A commonly used feature map:
\[
U(x) = \prod_i e^{i x_i Z_i}
\prod_{i<j} e^{i x_i x_j Z_i Z_j}.
\]

This corresponds to:
\[
H(x) = \sum_i x_i Z_i + \sum_{i<j} x_i x_j Z_i Z_j.
\]

\begin{itemize}
\item Linear terms $\rightarrow$ local rotations
\item Quadratic terms $\rightarrow$ entangling interactions
\end{itemize}
\end{frame}

%------------------------------------------------

\begin{frame}{Connection to Many-Body Physics}
The feature map Hamiltonian resembles a spin model:
\[
H(x) = \sum_i h_i(x) Z_i + \sum_{ij} J_{ij}(x) Z_i Z_j.
\]

Interpretation:
\begin{itemize}
\item Data acts as coupling constants
\item Feature map = time evolution under data-dependent Hamiltonian
\end{itemize}
\end{frame}

%================================================
\section{Kernel Expressivity}
%================================================

\begin{frame}{Kernel Expressivity}
The kernel is
\[
K(x,x') = |\braket{\phi(x)}{\phi(x')}|^2.
\]

Expressivity refers to:
\begin{itemize}
\item how well the kernel separates data
\item richness of the induced feature space
\end{itemize}
\end{frame}

%------------------------------------------------

\begin{frame}{Expansion of the Kernel}
Using
\[
\ket{\phi(x)} = U(x)\ket{0},
\]
we have
\[
K(x,x') = |\bra{0} U^\dagger(x) U(x') \ket{0}|^2.
\]

This depends on the operator:
\[
U^\dagger(x)U(x') = e^{i(H(x') - H(x)) + \cdots}.
\]

Thus the kernel encodes differences in the effective Hamiltonians.
\end{frame}

%------------------------------------------------

\begin{frame}{Expressivity vs Circuit Depth}
\begin{itemize}
\item shallow circuits:
\begin{itemize}
\item limited correlations
\item simple kernels
\end{itemize}
\item deep circuits:
\begin{itemize}
\item highly nonlinear kernels
\item potentially hard to simulate classically
\end{itemize}
\end{itemize}
\end{frame}

%================================================
\section{Entanglement and Expressivity}
%================================================

\begin{frame}{Entanglement as a Resource}
Quantum feature maps can generate entanglement:
\[
\ket{\phi(x)} = U(x)\ket{0}.
\]

Measure using von Neumann entropy:
\[
S(\rho_A) = -\mathrm{Tr}(\rho_A \log \rho_A).
\]
\end{frame}

%------------------------------------------------

\begin{frame}{Expressivity vs Entanglement}
\begin{itemize}
\item No entanglement:
\begin{itemize}
\item kernel factorizes
\item limited expressivity
\end{itemize}
\item Moderate entanglement:
\begin{itemize}
\item richer correlations
\item improved classification power
\end{itemize}
\item Too much entanglement:
\begin{itemize}
\item concentration of measure
\item kernel becomes nearly constant
\end{itemize}
\end{itemize}
\end{frame}

%------------------------------------------------

\begin{frame}{Concentration of Measure}
For highly random states:
\[
|\braket{\phi(x)}{\phi(x')}|^2 \approx \frac{1}{2^n}.
\]

\begin{itemize}
\item kernel loses discrimination power
\item analogous to barren plateaus
\end{itemize}
\end{frame}

%================================================
\section{Quantum Fisher Information}
%================================================

\begin{frame}{Quantum Fisher Information (QFI)}
The QFI measures sensitivity of states:
\[
F_{ij}(x) = 4\,\mathrm{Re}
\left(
\langle \partial_i \phi | \partial_j \phi \rangle
-
\langle \partial_i \phi | \phi \rangle
\langle \phi | \partial_j \phi \rangle
\right).
\]
\end{frame}

%------------------------------------------------

\begin{frame}{Geometric Interpretation}
\begin{itemize}
\item QFI defines a metric on quantum states
\item determines distinguishability of feature states
\end{itemize}

\[
ds^2 = \sum_{ij} F_{ij} dx_i dx_j
\]
\end{frame}

%------------------------------------------------

\begin{frame}{QFI and Kernel}
For small perturbations:
\[
K(x,x+\delta x) \approx 1 - \frac{1}{4} \delta x^T F \delta x.
\]

Thus:
\begin{itemize}
\item QFI controls local kernel geometry
\item large QFI $\Rightarrow$ better discrimination
\end{itemize}
\end{frame}

%------------------------------------------------

\begin{frame}{Connection to Learning}
\begin{itemize}
\item QFI $\leftrightarrow$ curvature of feature space
\item determines resolution of classification boundary
\end{itemize}
\end{frame}

%================================================
\section{Comparison with Classical Kernels}
%================================================

\begin{frame}{Classical Kernels}
Common examples:
\begin{itemize}
\item Linear:
\[
K(x,x') = x \cdot x'
\]
\item Polynomial:
\[
K(x,x') = (x \cdot x' + c)^d
\]
\item RBF:
\[
K(x,x') = e^{-\|x-x'\|^2 / 2\sigma^2}
\]
\end{itemize}
\end{frame}

%------------------------------------------------

\begin{frame}{Quantum vs Classical Kernels}
\begin{tabular}{l|l|l}
 & Classical & Quantum \\
\hline
Feature space & polynomial/infinite & $2^n$ dimensional \\
Computation & explicit/implicit & quantum overlaps \\
Correlations & limited & entangled \\
\end{tabular}
\end{frame}

%------------------------------------------------

\begin{frame}{When Quantum Helps}
Potential advantage if:
\begin{itemize}
\item kernel is hard to compute classically
\item feature map generates non-classical correlations
\item data matches quantum structure
\end{itemize}
\end{frame}

%------------------------------------------------

\begin{frame}{Limitations}
\begin{itemize}
\item kernel estimation cost
\item noise
\item unclear advantage for generic datasets
\end{itemize}
\end{frame}

%================================================
\section{Summary of Advanced Concepts}
%================================================

\begin{frame}{Summary}
\begin{itemize}
\item Feature maps = data-dependent Hamiltonians
\item Kernel = overlap of quantum states
\item Expressivity tied to entanglement
\item QFI defines geometry of feature space
\item Quantum kernels can exceed classical ones in structure
\end{itemize}
\end{frame}

\end{document}

